\chapter{统一的 token-level Mirror Descent 与 OPD 实例化}
\label{ch:four-methods}

本章只使用一套序列建模符号。给定 prompt $x$，student 生成 rollout $y_{1:H}$；第 $t$ 个 token 的历史为 $s_t=(x,y_{<t})$。固定第 $k$ 个 outer step 的 anchor policy $\pi_k$，候选 student policy 记为 $\pi_\theta$，teacher policy 记为 $T$，镜像正则的温度参数记为 $\eta>0$。$\eta$ 决定 mirror target 离 anchor 和 teacher 的位置；MDPO 的参数更新步长另记为 $\alpha$。$s_t$ 只表示 rollout 的条件历史，训练样本始终是一条完整的 prompt--rollout；四种 inner 的区别只来自求解同一个 mirror 子问题的方式。

\section{共同的 mirror 子问题与精确解}
\label{sec:ch2-common}

\begin{definition}[Token-level mirror 子问题]
\label{def:ch2-mirror-subproblem}
给定固定历史 $s$、token 信号 $f(s,a)$ 和 anchor $\pi_k$，定义
\begin{equation}
\label{eq:ch2-mirror}
\mathcal M_s(p;f,\pi_k)
=\mathbb E_{a\sim p(\cdot\mid s)}[f(s,a)]
-\frac1\eta\operatorname{KL}\!\left(p(\cdot\mid s)\middle\Vert\pi_k(\cdot\mid s)\right).
\end{equation}
镜像子问题是在所有满足 $p\ll\pi_k$ 的概率分布上最大化 $\mathcal M_s$。同时定义配分函数
\begin{equation}
\label{eq:ch2-mirror-normalizer}
Z_f(s)=\sum_{v\in\mathcal V}\pi_k(v\mid s)\exp\!\bigl(\eta f(s,v)\bigr).
\end{equation}
\end{definition}

\begin{theorem}[KL mirror 子问题的 Gibbs 解]
\label{thm:ch2-gibbs}
在有限词表、$Z_f(s)<\infty$ 且 $\pi_k(v\mid s)>0$ 的条件下，式 \eqref{eq:ch2-mirror} 的唯一解是
\begin{equation}
\label{eq:ch2-optimal-policy}
\pi_k^+(v\mid s)=\frac{\pi_k(v\mid s)e^{\eta f(s,v)}}{Z_f(s)},
\qquad
\log\frac{\pi_k^+(v\mid s)}{\pi_k(v\mid s)}
=\eta f(s,v)-\log Z_f(s).
\end{equation}
对任意候选 $p$，还有精确分解
\begin{equation}
\label{eq:ch2-kl-decomp}
\mathcal M_s(p;f,\pi_k)
=\frac1\eta\log Z_f(s)
-\frac1\eta\operatorname{KL}\!\left(p\middle\Vert\pi_k^+(\cdot\mid s)\right).
\end{equation}
\end{theorem}

\begin{proof}
由式 \eqref{eq:ch2-optimal-policy}，有
\[
\log\frac{p(v\mid s)}{\pi_k^+(v\mid s)}
=\log\frac{p(v\mid s)}{\pi_k(v\mid s)}-\eta f(s,v)+\log Z_f(s).
\]
对 $v$ 按 $p$ 求期望，得到
\[
\operatorname{KL}(p\Vert\pi_k^+)
=\operatorname{KL}(p\Vert\pi_k)-\eta\mathbb E_p[f]+\log Z_f(s).
\]
移项即得式 \eqref{eq:ch2-kl-decomp}。KL 非负且仅在 $p=\pi_k^+$ 时取零，因此最优解唯一。
\end{proof}

\begin{definition}[中心化 token 信号]
\label{def:ch2-centered-signal}
对任意 $f$ 定义 anchor 条件均值和中心化信号
\[
\bar f_k(s)=\mathbb E_{a\sim\pi_k}[f(s,a)\mid s],
\qquad
\widetilde f_k(s,a)=f(s,a)-\bar f_k(s).
\]
若 $f=A^{\pi_k}$ 是 anchor 的优势函数，则 $\bar f_k(s)=0$ 是优势函数定义推出的恒等式；对一般 teacher score 或 OPD log-ratio，$\bar f_k(s)=0$ 不是自动成立的假设。
\end{definition}

\begin{proposition}[优势函数的 anchor 条件均值]
\label{prop:ch2-advantage-centering}
若 $f(s,a)=A^{\pi_k}(s,a)$，则对每个满足 $\pi_k(\cdot\mid s)$ 定义的历史，
\begin{equation}
\label{eq:ch2-advantage-zero-mean}
\mathbb E_{a\sim\pi_k(\cdot\mid s)}[f(s,a)\mid s]=0.
\end{equation}
对于一般 OPD log-ratio 或 teacher score，该等式不是定义恒等式，必须显式估计并减去条件均值。
\end{proposition}

\begin{proof}
由 $A^{\pi_k}(s,a)=Q^{\pi_k}(s,a)-V^{\pi_k}(s)$ 以及
$V^{\pi_k}(s)=\mathbb E_{a\sim\pi_k}[Q^{\pi_k}(s,a)\mid s]$，有
\begin{align*}
\mathbb E_{\pi_k}[A^{\pi_k}(s,a)\mid s]
&=\mathbb E_{\pi_k}[Q^{\pi_k}(s,a)\mid s]-V^{\pi_k}(s)\\
&=V^{\pi_k}(s)-V^{\pi_k}(s)=0.
\end{align*}
\end{proof}

\begin{proposition}[中心化不改变 mirror 最优策略]
\label{prop:ch2-centering}
用 $\widetilde f_k$ 替代 $f$ 不改变式 \eqref{eq:ch2-mirror} 的最优分布。
\end{proposition}

\begin{proof}
因为 $\bar f_k(s)$ 与 action 无关，
\[
\pi_k(v\mid s)e^{\eta\widetilde f_k(s,v)}
=e^{-\eta\bar f_k(s)}\pi_k(v\mid s)e^{\eta f(s,v)}.
\]
分子和归一化常数同时乘以同一个因子 $e^{-\eta\bar f_k(s)}$，归一化后的分布不变。
\end{proof}

\section{从镜像最优性方程到四种可训练目标}
\label{sec:ch2-objective-derivation}

定理 \ref{thm:ch2-gibbs} 给出的不是一个抽象的分布结论，而是一个可以被监督学习近似的逐 token 方程。把式 \eqref{eq:ch2-optimal-policy} 移项，得到
\begin{equation}
\label{eq:ch2-logratio-equation}
\log\pi_k^+(a\mid s)-\log\pi_k(a\mid s)
=\eta f(s,a)-\log Z_f(s).
\end{equation}
左侧是 candidate student 相对于冻结 anchor 的 log-ratio，右侧只由冻结的 token signal 和归一化常数决定。

\begin{definition}[Mirror log-ratio target]
\label{def:ch2-logratio-target}
定义候选 student 相对于冻结 anchor 的 log-ratio 与 Gibbs target 为
\[
U_\theta(s,a)=\log\pi_\theta(a\mid s)-\log\pi_k(a\mid s),
\qquad
H_k(s,a)=\eta f(s,a)-\log Z_f(s).
\]
式 \eqref{eq:ch2-optimal-policy} 等价于 $U_\theta=H_k$；因此 $H_k$ 是参数化 inner 所要实现的逐 token 最优性方程。
\end{definition}

由这个定义，求解 mirror 子问题在参数化策略类中的直接替代，就是在训练测度下让 $U_\theta$ 拟合 $H_k$：
\begin{equation}
\label{eq:ch2-regression-from-mirror}
\min_\theta\ \mathbb E\left[(U_\theta-H_k)^2\right].
\end{equation}
因此 PMD-Part 的平方损失不是额外假设，而是把镜像最优性方程变成参数化模型可优化的投影问题。对式 \eqref{eq:ch2-regression-from-mirror} 求导，使用
\[
\nabla_\theta U_\theta(s,a)=\nabla_\theta\log\pi_\theta(a\mid s),
\]
立即得到式 \eqref{eq:ch2-part-gradient}。

PMD-Mean 来自同一个方程的中心化。由于 $\log Z_f(s)$ 只依赖 $s$，对 anchor action 取条件均值时
\[
\mathbb E_{\pi_k}[H_k(s,a)\mid s]
=\eta\mathbb E_{\pi_k}[f(s,a)\mid s]-\log Z_f(s).
\]
减去该均值后，未知的 $\log Z_f$ 消失，得到
\begin{equation}
\label{eq:ch2-centered-logratio-equation}
H_k(s,a)-\mathbb E_{\pi_k}[H_k(s,\cdot)\mid s]
=\eta\widetilde f_k(s,a).
\end{equation}
这解释了 PMD-Mean 的 target 来源。它不能只回归中心化 target，还必须保留 $U_\theta=\log(\pi_\theta/\pi_k)$ 的概率归一化约束；否则相差一个依赖历史的常数的所有 log-ratio 都会被视为同一个解。

BPO 再对式 \eqref{eq:ch2-logratio-equation} 两边作同样的中心化，但中心化的是 residual：
\begin{align}
D_\theta(s,a)
&=\bigl[\eta f(s,a)-U_\theta(s,a)\bigr]
-\mathbb E_{\pi_k}\bigl[\eta f(s,\cdot)-U_\theta(s,\cdot)\mid s\bigr].
\label{eq:ch2-bpo-from-optimality}
\end{align}
当 $D_\theta=0$ 时，只能先推出 $U_\theta=\eta f-C(s)$；概率归一化再唯一确定 $C(s)=\log Z_f(s)$。这就是 BPO 与 PMD-Part 共享精确零点、但有限参数更新不同的原因。BPO 的 score centering 并非预先加入的 baseline，而是对 $\mathbb E_{\pi_k}[U_\theta]$ 求导后必然产生的项。

MDPO 则不把式 \eqref{eq:ch2-logratio-equation} 改写成回归，而直接对原始 mirror objective 求最大化。由式 \eqref{eq:ch2-kl-decomp}，对任意参数化 student 有
\begin{equation}
\label{eq:ch2-mdpo-kl-gap}
\mathcal J_{\mathrm{MDPO}}(\theta)
=\frac1\eta\log Z_f(s)
-\frac1\eta\operatorname{KL}\!\left(\pi_\theta(\cdot\mid s)\middle\Vert\pi_k^+(\cdot\mid s)\right).
\end{equation}
所以 MDPO 的 inner 是在直接最小化 student 与 Gibbs target 的 reverse KL gap；PMD-Part、PMD-Mean 和 BPO 则是在不同坐标和不同投影测度下最小化该最优性方程的平方残差。

\begin{proposition}[平方回归的投影含义]
\label{prop:ch2-regression-projection}
若 student 函数类不能表示 $H_k$，PMD-Part 的精确 inner 解是 $H_k$ 在函数类诱导的 $L_2(d_{\pi_k}\pi_k)$ 几何下的投影；若训练样本的历史访问测度改为 $m$，投影也相应变为 $L_2(m\pi_k)$。因此 inner loss 为零只在 target 可表示且训练测度覆盖目标测度时才意味着精确 mirror 解。
\end{proposition}

\begin{proof}
平方损失是 Hilbert 空间中的正交投影问题。对任意可行扰动 $v$，最优点 $U^\star$ 满足
\[
\mathbb E_{m,\pi_k}[(U^\star-H_k)v]=0.
\]
这正是 $H_k$ 到可行函数集合的正交投影条件。若 $H_k$ 属于该集合，投影残差为零；否则残差由函数类和测度共同决定。
\end{proof}

\section{四种 inner 的统一推导}
\label{sec:ch2-four-methods}

下面固定一个历史，令
\[
\ell^T(a)=\log T(a\mid s),\qquad
\ell^k(a)=\log\pi_k(a\mid s),\qquad
\ell^\theta(a)=\log\pi_\theta(a\mid s).
\]
训练时只对 $\ell^\theta$ 反向传播；teacher、anchor、目标常数和所有估计的 KL 都停止梯度。

\subsection{PMD-Part：回归完整 Gibbs log-ratio}
\label{sec:ch2-pmd-part}

定义显式 target
\[
H_k(s,a)=\eta f(s,a)-\log Z_f(s),
\qquad
U_\theta(s,a)=\log\pi_\theta(a\mid s)-\log\pi_k(a\mid s).
\]
PMD-Part 在 anchor 测度下最小化
\begin{equation}
\label{eq:ch2-part-loss}
\mathcal L_{\mathrm{Part}}(\theta)
=\frac12\mathbb E_{s\sim d_k,\,a\sim\pi_k}
\left[U_\theta(s,a)-H_k(s,a)\right]^2.
\end{equation}
对 target 停止梯度，链式法则给出
\begin{equation}
\label{eq:ch2-part-gradient}
-\nabla_\theta\mathcal L_{\mathrm{Part}}
=\mathbb E_{d_k,\pi_k}
\left[H_k(s,a)-U_\theta(s,a)\right]
\nabla_\theta\log\pi_\theta(a\mid s).
\end{equation}
这一步是普通的 log-ratio regression；它没有把 KL 视为常数，也没有使用 action-level importance ratio。

固定 history 后令 $q_\theta(a)=U_\theta(s,a)-H_k(s,a)$。由于 anchor 测度与 target 都冻结，逐步求导得到
\begin{align}
\nabla_\theta\mathcal L_{\mathrm{Part}}(s)
&=\frac12\nabla_\theta\sum_a\pi_k(a\mid s)q_\theta(a)^2\\
&=\sum_a\pi_k(a\mid s)q_\theta(a)\nabla_\theta U_\theta(s,a)\\
&=\E_{a\sim\pi_k}[q_\theta(a)\nabla_\theta\log\pi_\theta(a\mid s)\mid s].
\end{align}
若实际 rollout 来自行为测度 $\beta_k$，而目标仍是 $L_2(\pi_k)$ 投影，则每个 sampled residual 乘以冻结权重 $\operatorname{sg}(\pi_k(a\mid s)/\beta_k(a\mid s))$。不加这个权重时，算法仍然有清晰含义，但求解的是 $L_2(\beta_k)$ 投影。

\begin{algorithm}[ht]
\caption{PMD-Part inner update\\{\small\textit{Semantics: 固定 anchor 下的完整 Gibbs log-ratio 回归}}}
\begin{algorithmic}[1]
\Require frozen $\pi_k$, teacher signal $f$, histories $\{s_i\}_{i=1}^B$, $\eta$, step size $\alpha$
\Ensure updated parameters $\theta$
\For{$i=1,\ldots,B$}
    \State sample $a_i\sim\pi_k(\cdot\mid s_i)$
    \State $\log Z_i\gets\operatorname{logsumexp}_{v\in\mathcal V}(\log\pi_k(v\mid s_i)+\eta f(s_i,v))$
    \State $H_i\gets\eta f(s_i,a_i)-\operatorname{stopgrad}(\log Z_i)$
    \State $U_i\gets\log\pi_\theta(a_i\mid s_i)-\operatorname{stopgrad}(\log\pi_k(a_i\mid s_i))$
    \State $\ell_i\gets\frac12\bigl(U_i-\operatorname{stopgrad}(H_i)\bigr)^2$
\EndFor
\State $\theta\gets\theta-\alpha\nabla_\theta(B^{-1}\sum_i\ell_i)$
\State \Return $\theta$
\end{algorithmic}
\end{algorithm}

\subsection{PMD-Mean：中心化 target 与归一化约束}
\label{sec:ch2-pmd-mean}

由定义可得
\[
\mathbb E_{a\sim\pi_k}[f(s,a)]=\bar f_k(s),
\qquad
\mathbb E_{a\sim\pi_k}[\widetilde f_k(s,a)]=0.
\]
PMD-Mean 用 $\eta\widetilde f_k$ 作为回归 target，并要求 student log-ratio 满足归一化约束
\begin{equation}
\label{eq:ch2-mean-loss}
\min_{U}\ \frac12\mathbb E_{\pi_k}\left[(U-\eta\widetilde f_k)^2\mid s\right]
\quad\text{s.t.}\quad
\mathbb E_{\pi_k}[e^{U}\mid s]=1.
\end{equation}

\begin{theorem}[PMD-Mean 的逐历史 Lambert-$W$ 投影]
\label{thm:ch2-mean-lambert}
在有限词表、$\pi_k(a\mid s)>0$ 且约束可行时，上述中心化平方回归的唯一 KKT 解满足
\begin{equation}
\label{eq:ch2-mean-lambert-theorem}
U_a=G_a-W\!\left(\lambda(s)e^{G_a}\right),\qquad
G_a=\eta\widetilde f_k(s,a),
\end{equation}
其中 $\lambda(s)\geq0$ 由 $\sum_a\pi_k(a\mid s)e^{U_a}=1$ 唯一确定。
\end{theorem}

\begin{proof}
以下给出固定历史 $s$ 的 KKT 推导。令词表上的函数值为 $U_a$ 和 $G_a=\eta\widetilde f_k(s,a)$，拉格朗日函数为
\[
\mathscr L(U,\lambda)
=\frac12\sum_a\pi_k(a\mid s)(U_a-G_a)^2
+\lambda\left(\sum_a\pi_k(a\mid s)e^{U_a}-1\right).
\]
一阶条件为
\[
U_a-G_a+\lambda e^{U_a}=0.
\]
令 $V_a=G_a-U_a$，则 $V_a e^{V_a}=\lambda e^{G_a}$，所以
\begin{equation}
\label{eq:ch2-mean-lambert}
U_a=G_a-W\!\left(\lambda e^{G_a}\right),
\end{equation}
其中 $\lambda$ 由归一化约束唯一确定。为说明唯一性，定义
\[
\Phi(\lambda)=\sum_a\pi_k(a\mid s)\exp\!\left(G_a-W(\lambda e^{G_a})\right),qquad \lambda\geq0.
\]
Lambert--$W$ 在非负轴上单调递增，因此每一项随 $\lambda$ 单调递减；当 $\lambda=0$ 时 $\Phi(0)=\sum_a\pi_k(a\mid s)e^{G_a}$，而 $\lambda\to\infty$ 时 $\Phi(\lambda)\to0$。只要中心化 target 不已经满足 $\Phi(0)=1$，方程 $\Phi(\lambda)=1$ 有唯一正根；若 $\Phi(0)=1$，则 $\lambda=0$，得到 $U_a=G_a$。因此 PMD-Mean 是中心化 target 在指数归一化集合上的投影。严格凸性和 $\Phi$ 的严格单调性给出 KKT 解的唯一性。
\end{proof}

约束 $\sum_a\pi_\theta(a\mid s)=1$ 正是 log-ratio 坐标中的指数归一化条件。特别地，精确 Gibbs 解的 log-ratio 为
\[
U_a^\star=\eta f(s,a)-\log Z_f(s),\qquad
U_a^\star-\eta\widetilde f_k(s,a)=\eta\bar f_k(s)-\log Z_f(s),
\]
它与 action 无关，且自动满足 $\sum_a\pi_k(a\mid s)e^{U_a^\star}=1$；这说明 PMD-Mean 恢复的是同一个 Gibbs 分布，而不是把 $\log Z_f$ 任意设为零。

实际样本中需要分别估计 $\bar U_\theta(s)$ 和 $\bar f_k(s)$。为避免一个 action 同时出现在 residual 与 center 中，给定 $K\ge2$ 个独立 anchor actions，使用 leave-one-out 均值
\[
\widehat{\bar U}_{-j}=\frac1{K-1}\sum_{\ell\ne j}U_\theta(s,a_\ell),\qquad
\widehat{\bar f}_{-j}=\frac1{K-1}\sum_{\ell\ne j}f(s,a_\ell).
\]
若 $K=1$，center 不可估计，应该退化为 PMD-Part 或使用上一批冻结的 center。center 中的 $U_\theta$ 必须保留梯度；把它整体 stop-gradient 会丢掉归一化约束对应的 centered-score 项。

\begin{algorithm}[ht]
\caption{PMD-Mean inner update\\{\small\textit{Semantics: 中心化 target 在指数归一化集合上的投影}}}
\begin{algorithmic}[1]
\Require frozen $\pi_k$, signal $f$, histories $\{s_i\}$, $K\ge2$, $\eta$, step size $\alpha$
\Ensure updated parameters $\theta$
\For{each history $s_i$}
    \For{$j=1,\ldots,K$}
        \State sample $a_{ij}\sim\pi_k(\cdot\mid s_i)$
        \State $U_{ij}\gets\log\pi_\theta(a_{ij}\mid s_i)-\operatorname{stopgrad}(\log\pi_k(a_{ij}\mid s_i))$
        \State $F_{ij}\gets f(s_i,a_{ij})$
    \EndFor
    \For{$j=1,\ldots,K$}
        \State $\bar U_{-j}\gets(K-1)^{-1}\sum_{\ell\ne j}U_{i\ell}$, $\bar F_{-j}\gets(K-1)^{-1}\sum_{\ell\ne j}F_{i\ell}$
        \State $\ell_{ij}\gets\frac12[(U_{ij}-\bar U_{-j})-\eta(F_{ij}-\bar F_{-j})]^2$
    \EndFor
\EndFor
\State $\theta\gets\theta-\alpha\nabla_\theta(BK)^{-1}\sum_{i,j}\ell_{ij}$
\State \Return $\theta$
\end{algorithmic}
\end{algorithm}


\subsection{BPO：中心化最优性 residual}
\label{sec:ch2-bpo}

BPO 不显式回归 $\log Z_f$，而将同一个最优性方程的 action-independent 部分消去。定义
\begin{equation}
\label{eq:ch2-bpo-residual}
D_\theta(s,a)
=\eta f(s,a)-U_\theta(s,a)
-\mathbb E_{v\sim\pi_k}[\eta f(s,v)-U_\theta(s,v)\mid s].
\end{equation}
直接展开可得
\[
D_\theta(s,a)
=\eta\widetilde f_k(s,a)-U_\theta(s,a)
+\mathbb E_{v\sim\pi_k}[U_\theta(s,v)\mid s].
\]
未线性化 BPO 的损失为
\begin{equation}
\label{eq:ch2-bpo-loss}
\mathcal L_{\mathrm{BPO}}(\theta)
=\frac1{2\eta}\mathbb E_{d_k,\pi_k}[D_\theta(s,a)^2].
\end{equation}
由于
\[
\nabla_\theta\mathbb E_{v\sim\pi_k}[U_\theta(s,v)]
=\mathbb E_{v\sim\pi_k}[\nabla_\theta\log\pi_\theta(v\mid s)],
\]
所以
\begin{equation}
\label{eq:ch2-bpo-gradient}
-\eta\nabla_\theta\mathcal L_{\mathrm{BPO}}
=\mathbb E_{d_k,\pi_k}\left[D_\theta(s,a)
\left(\nabla_\theta\log\pi_\theta(a\mid s)
-\mathbb E_{v\sim\pi_k}[\nabla_\theta\log\pi_\theta(v\mid s)\mid s]\right)\right].
\end{equation}
BPO 的 score centering 来自对当前 student--anchor KL 的精确求导；它不是因为先对 reward 做了一个近似。

令 $G_\theta(s)=\E_{\pi_k}[\eta f(s,v)-U_\theta(s,v)\mid s]$，则
\[
\nabla_\theta D_\theta(s,a)=-\nabla_\theta U_\theta(s,a)+\nabla_\theta\E_{\pi_k}[U_\theta(s,v)\mid s]
=-g_\theta(s,a)+\bar g_\theta(s).
\]
所以 centered score 是对 center 本身求导的结果，而不是把 reward 任意减去 baseline。无偏的 Monte Carlo 实现最好为每个 history 使用独立的 action $a_i$ 计算 residual，再使用独立集合 $\{b_{i\ell}\}$ 估计 $G_\theta(s_i)$；若复用样本，必须用 leave-one-out center。

\begin{algorithm}[ht]
\caption{BPO inner update\\{\small\textit{Semantics: 对中心化最优性 residual 求平方损失}}}
\begin{algorithmic}[1]
\Require frozen $\pi_k$, signal $f$, histories $\{s_i\}$, center size $K$, $\eta$, step size $\alpha$
\Ensure updated parameters $\theta$
\For{each history $s_i$}
    \State sample $a_i\sim\pi_k(\cdot\mid s_i)$ and $b_{i1},\ldots,b_{iK}\sim\pi_k(\cdot\mid s_i)$ independently
    \State $z_i\gets\eta f(s_i,a_i)-\log\pi_\theta(a_i\mid s_i)+\operatorname{stopgrad}(\log\pi_k(a_i\mid s_i))$
    \State $\bar z_i\gets K^{-1}\sum_{\ell=1}^K[\eta f(s_i,b_{i\ell})-\log\pi_\theta(b_{i\ell}\mid s_i)+\operatorname{stopgrad}(\log\pi_k(b_{i\ell}\mid s_i))]$
    \State $d_i\gets z_i-\bar z_i$, $\ell_i\gets d_i^2/(2\eta)$
\EndFor
\State $\theta\gets\theta-\alpha\nabla_\theta B^{-1}\sum_i\ell_i$
\State \Return $\theta$
\end{algorithmic}
\end{algorithm}

若从行为测度 $\beta$ 采样，$z_i$ 和每个 center summand 都乘以冻结的 $\operatorname{sg}(\pi_k/\beta)$ 才对应 $\pi_k$ 期望。该权重只改变积分测度，不能对它反向传播。

\subsection{MDPO：直接优化 mirror objective}
\label{sec:ch2-mdpo-generic}

MDPO 直接最大化式 \eqref{eq:ch2-mirror} 的参数化版本：
\begin{equation}
\label{eq:ch2-mdpo-generic-objective}
\mathcal J_{\mathrm{MDPO}}(\theta)
=\mathbb E_{a\sim\pi_\theta}[f(s,a)]
-\frac1\eta\operatorname{KL}\!\left(\pi_\theta(\cdot\mid s)\middle\Vert\pi_k(\cdot\mid s)\right).
\end{equation}
完整词表时 KL 是
\begin{equation}
\label{eq:ch2-mdpo-full-kl}
\operatorname{KL}(\pi_\theta\Vert\pi_k)
=\sum_a\pi_\theta(a\mid s)
\left[\log\pi_\theta(a\mid s)-\log\pi_k(a\mid s)\right].
\end{equation}
若 action 从 anchor 采样，定义 $R_\theta(s,a)=\exp(\ell^\theta(a)-\ell^k(a))$，则 action-level IS 形式为
\begin{equation}
\label{eq:ch2-mdpo-sample-integrand}
\widehat{\mathcal J}_{\mathrm{MDPO}}
=\mathbb E_{a\sim\pi_k}\left[R_\theta(s,a)f(s,a)-\frac1\eta R_\theta(s,a)\log R_\theta(s,a)\right].
\end{equation}
其 ascent direction 为
\begin{equation}
\label{eq:ch2-mdpo-generic-gradient}
\nabla_\theta\mathcal J_{\mathrm{MDPO}}
=\mathbb E_{\pi_k}\left[R_\theta(s,a)\left(f(s,a)-\frac1\eta\log R_\theta(s,a)\right)\nabla_\theta\log\pi_\theta(a\mid s)\right].
\end{equation}
额外的 $-R_\theta/\eta$ 项乘以 score 后，在完整 action 求和中因 $\mathbb E_{\pi_\theta}[\nabla\log\pi_\theta]=0$ 消失。MDPO 因而直接优化 mirror objective，其他三种方法优化的是由同一最优性方程构造的 regression objective。

MDPO 的自然梯度只改变更新几何，不改变上面的 token target。令 $g=\nabla_\theta\mathcal J_{\mathrm{MDPO}}$，并令
\[
F_\theta=\E_{a\sim\pi_\theta}[g_\theta(a)g_\theta(a)^\top],\qquad
g_\theta(a)=\nabla_\theta\log\pi_\theta(a\mid s).
\]
若从 anchor 采样，用 $\operatorname{sg}(R_\theta)$ 进行测度校正。由于 categorical Fisher 的零空间是 logits 的常数方向，$F_\theta^\dagger g$ 位于 simplex 切空间；神经网络参数化下则需用阻尼逆 $(F+\lambda I)^{-1}$ 或 KL trust region，不能把自然步长与 mirror 温度 $\eta$ 混用。

\begin{algorithm}[ht]
\caption{MDPO natural-gradient inner update\\{\small\textit{Semantics: 直接优化固定 history 的 mirror objective}}}
\begin{algorithmic}[1]
\Require frozen $\pi_k$, signal $f$, histories $\{s_i\}$, $\eta$, natural step $\alpha$
\Ensure updated parameters $\theta$
\For{each history $s_i$}
    \State sample $a_i\sim\pi_k(\cdot\mid s_i)$
    \State $R_i\gets\exp(\log\pi_\theta(a_i\mid s_i)-\operatorname{stopgrad}(\log\pi_k(a_i\mid s_i)))$
    \State $j_i\gets R_i f(s_i,a_i)-\eta^{-1}R_i\log R_i$ \Comment{保留 $R_i$ 对 student 的梯度}
\EndFor
\State $g\gets\nabla_\theta(B^{-1}\sum_i j_i)$
\State $g_i^{\mathrm{score}}\gets\nabla_\theta\log\pi_\theta(a_i\mid s_i)$
\State $F\gets B^{-1}\sum_i\operatorname{stopgrad}(R_i)g_i^{\mathrm{score}}(g_i^{\mathrm{score}})^\top$
\State $\theta\gets\theta+\alpha(F+\lambda I)^\dagger g$
\State \Return $\theta$
\end{algorithmic}
\end{algorithm}

\begin{theorem}[四种 inner 的共同精确解]
\label{thm:ch2-inner-common-root}
在 tabular policy、完整 action 支持和精确词表期望下，PMD-Part、PMD-Mean、未线性化 BPO 和 MDPO 都以式 \eqref{eq:ch2-optimal-policy} 为精确解。PMD-Mean 的结论要求其投影 target 可实现；参数化 student 或有限样本时，四者得到的是不同的投影或优化误差。
\end{theorem}

\begin{proof}
PMD-Part 在零损失处满足 $U_\theta=\eta f-\log Z_f$，于是得到式 \eqref{eq:ch2-optimal-policy}。BPO 的零 residual 只确定 $U_\theta=\eta f-C(s)$；代入 $\mathbb E_{\pi_k}[e^{U_\theta}]=1$ 得 $C(s)=\log Z_f(s)$。PMD-Mean 若中心化 target 属于归一化约束集合，则其平方投影残差为零，并由中心化不变性得到同一策略。MDPO 由式 \eqref{eq:ch2-kl-decomp} 直接得到唯一最大点。
\end{proof}


\begin{proposition}[未线性化 BPO 与 PMD-Part 的梯度差异]
\label{prop:ch2-bpo-part-gradient-gap}
固定一个历史 $s$，令
\begin{equation}
\label{eq:ch2-bpo-part-qc}
q(s,a)=\eta f(s,a)-\log Z_f(s)-U_\theta(s,a),
\qquad
c(s)=\log Z_f(s)+\mathbb E_{\pi_k}[U_\theta(s,\cdot)\mid s]-\eta\mathbb E_{\pi_k}[f(s,\cdot)\mid s].
\end{equation}
则未线性化 BPO residual 满足 $D_\theta=q+c$，并且
\begin{equation}
\label{eq:ch2-bpo-part-gradient-gap}
-\eta\nabla_\theta\mathcal L_{\mathrm{BPO}}
=-\nabla_\theta\mathcal L_{\mathrm{Part}}
+c(s)\,\mathbb E_{a\sim\pi_k}[\nabla_\theta\log\pi_\theta(a\mid s)\mid s].
\end{equation}
因此两者在精确镜像解处都有零梯度，但在一般参数点不相同。它们只有在
$c(s)\,\mathbb E_{\pi_k}[\nabla_\theta\log\pi_\theta]=0$
时才给出相同的方向。
\end{proposition}

\begin{proof}
从 PMD-Part 的最优性残差出发，先把 BPO 的中心化式逐项展开：
\begin{align*}
D_\theta(s,a)
&=\eta f(s,a)-U_\theta(s,a)
 -\mathbb E_{\pi_k}[\eta f(s,\cdot)-U_\theta(s,\cdot)\mid s]\\
&=\eta f(s,a)-\log Z_f(s)-U_\theta(s,a)
 +\log Z_f(s)+\mathbb E_{\pi_k}[U_\theta(s,\cdot)\mid s]\\
&=q(s,a)+c(s).
\end{align*}
由定义，$\mathbb E_{\pi_k}[D_\theta\mid s]=0$。BPO 的精确梯度为
\begin{align*}
-\eta\nabla_\theta\mathcal L_{\mathrm{BPO}}
&=\mathbb E_{\pi_k}\left[D_\theta
 \left(g_\theta-\bar g_\theta\right)\mid s\right],\\
 g_\theta(a)&=\nabla_\theta\log\pi_\theta(a\mid s),
 \qquad \bar g_\theta=\mathbb E_{\pi_k}[g_\theta\mid s].
\end{align*}
代入 $D_\theta=q+c$，并使用 $\mathbb E_{\pi_k}[D_\theta\mid s]=0$，可得
\begin{align*}
\mathbb E_{\pi_k}[D_\theta(g_\theta-\bar g_\theta)\mid s]
&=\mathbb E_{\pi_k}[(q+c)g_\theta\mid s]
 -\mathbb E_{\pi_k}[D_\theta\mid s]\bar g_\theta\\
&=\mathbb E_{\pi_k}[qg_\theta\mid s]+c(s)\bar g_\theta.
\end{align*}
而 PMD-Part 的梯度式为
$-\nabla_\theta\mathcal L_{\mathrm{Part}}=\mathbb E_{\pi_k}[qg_\theta\mid s]$，故式
\eqref{eq:ch2-bpo-part-gradient-gap}
成立。若 $q=0$，则由归一化约束有 $c=0$，两边同时为零；离开该点时 $c$ 通常非零，差异项保留。
\end{proof}

这个命题说明“都在 token level”只保证它们面对同一个局部镜像根，不保证优化轨迹相同。PMD-Part 回归完整的 action-dependent log-ratio，并把 $\log Z_f$ 当作固定 target；BPO 先消去 action-independent 常数，再对当前 student 的归一化项求导，因此多出 centered score。若使用的参数化使
$\mathbb E_{\pi_k}[\nabla_\theta\log\pi_\theta]=0$
（例如采样测度与当前 student 一致，或在特定 tabular 点上），差异项才会消失。

\section{OPD 特化：逐 token reward 与两种信息访问模式}
\label{sec:ch2-opd-specialization}

在 OPD 中，第 $k$ 个 outer step 先固定 anchor $\pi_k$，然后从它产生 rollout。对 prefix
$s_t=(x,y_{<t})$ 和实际 sampled token $a_t=y_t$，最直接的局部 teacher reward 是
\begin{equation}
\label{eq:ch2-opd-raw-reward}
r_k(s_t,a_t)=\log T(a_t\mid s_t)-\log\pi_k(a_t\mid s_t).
\end{equation}
它满足序列分解
\[
\sum_{t=1}^H r_k(s_t,a_t)=\log T(y\mid x)-\log\pi_k(y\mid x).
\]
因此在共同 mirror 子问题中应取 $f=r_k$；不能把整条序列的 log-ratio 当作每一个 prefix 的同一个 reward，否则会重复计算后缀差异。

可以对 reward 做两种不同的“归一化”，必须区分它们的理论作用。对任意只依赖 prefix 的 baseline $b_k(s)$，令
\begin{equation}
\label{eq:ch2-opd-centered-reward}
f_k^{\mathrm{ctr}}(s,a)=r_k(s,a)-b_k(s).
\end{equation}
它只把 partition function 乘上 $e^{-\eta b_k(s)}$，所以 Gibbs policy 不变；取
$b_k(s)=\mathbb E_{\pi_k}[r_k(s,\cdot)\mid s]$ 就是条件中心化。若再除以尺度 $\sigma_k(s)>0$，则
\[
f_k^{\mathrm{std}}=(r_k-b_k)/\sigma_k
\]
会把有效温度从 $\eta$ 改成 $\eta/\sigma_k$，不再是同一个 mirror step，除非同时把正则参数改为 $\eta\sigma_k$。因此 OPD 的默认实现只做 additive centering；标准化只能作为明确的 temperature schedule。

\subsection{完整 vocabulary logits：精确 token mirror target}
\label{sec:ch2-opd-full-vocab}

如果 teacher 和 anchor 都能在 prefix $s$ 返回完整词表 logits，则对每个 $v\in\mathcal V$ 计算
\[
r_k(s,v)=\log T(v\mid s)-\log\pi_k(v\mid s),
\qquad
\log Z_k(s)=\operatorname{logsumexp}_{v\in\mathcal V}
\bigl((1-\eta)\log\pi_k(v\mid s)+\eta\log T(v\mid s)\bigr).
\]
于是 exact target 为
\[
\pi_k^+(v\mid s)=\exp\bigl((1-\eta)\log\pi_k(v\mid s)+\eta\log T(v\mid s)-\log Z_k(s)\bigr).
\]
这时 PMD-Part 的 target、MDPO 的 reverse KL、以及 BPO/PMD-Mean 所需的 anchor conditional center 都可以在同一个 prefix 上精确计算。完整词表路径不需要 action-level importance sampling；唯一的随机性来自 prefix batch 和 optimizer。

在写算法之前先把四个 inner 的计算量化简。对一个固定 prefix，teacher、anchor 和当前 student 的信息都直接写成 token log-probability。PMD 的归一化常数是
\begin{equation}
\label{eq:ch2-opd-full-logz}
\log Z_k(s)=
\log\left(\sum_{v\in\mathcal V}
\exp\left((1-\eta)\log\pi_k(v\mid s)+\eta\log T(v\mid s)\right)\right).
\end{equation}
因此 exact mirror target 为
\[
\pi_k^+(v\mid s)=
\exp\!\left((1-\eta)\log\pi_k(v\mid s)+\eta\log T(v\mid s)-\log Z_k(s)\right).
\]
为了写出 PMD-Mean 和 BPO 的中心化形式，先把两个均值写成 KL 散度。定义
\begin{align}
\mathcal K_k^T(s)
&\coloneqq \operatorname{KL}\!\left(\pi_k(\cdot\mid s)\middle\Vert T(\cdot\mid s)\right),\\
\mathcal K_k^\theta(s)
&\coloneqq \operatorname{KL}\!\left(\pi_k(\cdot\mid s)\middle\Vert\pi_\theta(\cdot\mid s)\right).
\end{align}
在 outer--inner 记号 $\theta_{k,n}$ 中，$k$ 是 outer step 编号，$n$
是该 outer step 内的 inner optimizer step 编号。因此
$\mathcal K_{k,n}^\theta(s)$ 简写
$\operatorname{KL}(\pi_k(\cdot\mid s)\Vert\pi_{\theta_{k,n}}(\cdot\mid s))$：
左侧分布 $\pi_k$ 在整个 inner loop 固定，右侧 student 随 $n$ 改变，故每次
inner 更新后都要用新参数重算这个量。相反，$\mathcal K_k^T(s)$ 在固定
outer step 内不随 $n$ 改变。这里的 $\mathcal K$ 表示 KL 散度；它和
Top-$K$ 估计器中的整数 $K$ 不是同一个对象。
于是
\begin{align}
\overline{\log(T/\pi_k)}(s)&=-\mathcal K_k^T(s),\\
\overline{\log(\pi_\theta/\pi_k)}(s)&=-\mathcal K_k^\theta(s).
\end{align}
后续可以对这两个 KL 分别采用 full vocabulary、top-$K$ 或 sampled-token 估计，但不能混淆它们与 MDPO 中 $\operatorname{KL}(\pi_\theta\Vert\pi_k)$ 的方向。
于是四个 loss 可直接化为
\begin{align}
\mathcal L_{\mathrm{Part}}(s)
&=\frac12\mathbb E_{v\sim\pi_k(\cdot\mid s)}
\left[\log\frac{\pi_\theta(v\mid s)}{\pi_k(v\mid s)}
-\eta\log\frac{T(v\mid s)}{\pi_k(v\mid s)}
+\log Z_k(s)\right]^2,
\label{eq:ch2-opd-full-part-simplified}\\
\mathcal L_{\mathrm{Mean}}(s)
&=\frac12\mathbb E_{v\sim\pi_k(\cdot\mid s)}
\Bigg[
\left(\log\frac{\pi_\theta(v\mid s)}{\pi_k(v\mid s)}
+\mathcal K_k^\theta(s)\right)
\\[-2pt]
&\hspace{35mm}
-\eta\left(\log\frac{T(v\mid s)}{\pi_k(v\mid s)}
+\mathcal K_k^T(s)\right)
\Bigg]^2,
\label{eq:ch2-opd-full-mean-simplified}\\
D_\theta(s,v)
&=\eta\left(\log\frac{T(v\mid s)}{\pi_k(v\mid s)}
+\mathcal K_k^T(s)\right)
\\[-2pt]
&\quad-\left(\log\frac{\pi_\theta(v\mid s)}{\pi_k(v\mid s)}
+\mathcal K_k^\theta(s)\right),
\label{eq:ch2-opd-full-bpo-residual}\\
\mathcal L_{\mathrm{BPO}}(s)&=\frac1{2\eta}\mathbb E_{v\sim\pi_k(\cdot\mid s)}[D_\theta(s,v)^2],
\label{eq:ch2-opd-full-bpo-simplified}\\
\mathcal J_{\mathrm{MDPO}}(s)
&=\mathbb E_{v\sim\pi_\theta(\cdot\mid s)}
\left[\log\frac{T(v\mid s)}{\pi_k(v\mid s)}
-\frac1\eta\log\frac{\pi_\theta(v\mid s)}{\pi_k(v\mid s)}\right].
\label{eq:ch2-opd-full-mdpo-simplified}
\end{align}
由于 $\mathbb E_{v\sim\pi_k(\cdot\mid s)}[D_\theta(s,v)]=0$，BPO 的完整词表梯度进一步化为
\begin{equation}
\label{eq:ch2-opd-full-bpo-gradient-simplified}
-\eta\nabla_\theta\mathcal L_{\mathrm{BPO}}(s)
=\mathbb E_{v\sim\pi_k(\cdot\mid s)}
\left[D_\theta(s,v)\nabla_\theta\log\pi_\theta(v\mid s)\right].
\end{equation}
因此完整 vocabulary 路径无需显式计算 centered score 的第二项；它在精确求和中严格抵消。

\subsubsection*{对应的梯度与 stop-gradient 约定}
下面把式~\eqref{eq:ch2-opd-full-part-simplified}--\eqref{eq:ch2-opd-full-mdpo-simplified}
写成可直接实现的 policy-gradient 形式。固定 outer step $k$ 时，$\pi_k$、$T$、
$\log Z_k$ 和 $\mathcal K_k^T$ 都不依赖当前参数 $\theta$，实现中应使用
$\operatorname{stopgrad}$。相反，$\mathcal K_k^\theta$ 必须保留对
$\log\pi_\theta$ 的梯度。完整词表下
\begin{align}
\label{eq:ch2-opd-grad-ktheta}
\mathcal K_k^\theta(s)
&=\mathbb E_{v\sim\pi_k(\cdot\mid s)}
\left[\operatorname{stopgrad}\!\left(\log\pi_k(v\mid s)\right)
-\log\pi_\theta(v\mid s)\right],
\\
\nabla_\theta\mathcal K_k^\theta(s)
&
=-\mathbb E_{v\sim\pi_k(\cdot\mid s)}
\left[\nabla_\theta\log\pi_\theta(v\mid s)\right].
\end{align}
在下面的 policy-gradient 实现中，$A_{\mathrm{Part}}$、
$A_{\mathrm{Mean}}$、$D_\theta$ 以及 MDPO 的 log-ratio 权重都必须整体
使用 $\operatorname{stopgrad}$。这不是说它们在原始平方损失或原始
MDPO objective 中没有参数依赖，而是这些依赖已经在前面的解析求导中处理过；
直接实现 policy-gradient 时若不停止梯度，就会重复求导并产生错误的二阶项。

对 PMD-Part，残差对 $\theta$ 的导数就是
$\nabla_\theta\log\pi_\theta(v\mid s)$，故
\begin{align}
\nabla_\theta\mathcal L_{\mathrm{Part}}(s)
&=\mathbb E_{v\sim\pi_k(\cdot\mid s)}
\left[\operatorname{stopgrad}\!\left(A_{\mathrm{Part}}(s,v)\right)
\nabla_\theta\log\pi_\theta(v\mid s)\right],\\
A_{\mathrm{Part}}(s,v)
&=\log\frac{\pi_\theta(v\mid s)}{\pi_k(v\mid s)}
-\eta\log\frac{T(v\mid s)}{\pi_k(v\mid s)}+\operatorname{stopgrad}(\log Z_k(s)).
\end{align}
因此其梯度下降方向是一个以 $\pi_k$ 为测度、以 $-A_{\mathrm{Part}}$ 为权重的
score-function 更新。

对 PMD-Mean，令 $A_{\mathrm{Mean}}(s,v)$ 表示式~\eqref{eq:ch2-opd-full-mean-simplified}
中方括号内的量。由式~\eqref{eq:ch2-opd-grad-ktheta}，
\begin{equation}
\label{eq:ch2-opd-grad-mean-residual}
\nabla_\theta A_{\mathrm{Mean}}(s,v)
=\nabla_\theta\log\pi_\theta(v\mid s)
-\mathbb E_{u\sim\pi_k(\cdot\mid s)}
\left[\nabla_\theta\log\pi_\theta(u\mid s)\right].
\end{equation}
于是
\begin{align}
\nabla_\theta\mathcal L_{\mathrm{Mean}}(s)
&=\mathbb E_{v\sim\pi_k(\cdot\mid s)}\Bigg[A_{\mathrm{Mean}}(s,v)\Bigg(
\nabla_\theta\log\pi_\theta(v\mid s)
-\mathbb E_{u\sim\pi_k(\cdot\mid s)}
\left[\nabla_\theta\log\pi_\theta(u\mid s)\right]\Bigg)\Bigg]\\
&=\mathbb E_{v\sim\pi_k(\cdot\mid s)}
\left[\operatorname{stopgrad}\!\left(A_{\mathrm{Mean}}(s,v)\right)
\nabla_\theta\log\pi_\theta(v\mid s)\right],
\end{align}
其中第二个等号使用了
$\mathbb E_{v\sim\pi_k}[A_{\mathrm{Mean}}(s,v)]=0$。这一步只有在
$\mathcal K_k^\theta$ 保留梯度时才成立。

因为 $D_\theta=-A_{\mathrm{Mean}}$，BPO 的梯度为
\begin{align}
\nabla_\theta\mathcal L_{\mathrm{BPO}}(s)
&=-\frac1\eta\mathbb E_{v\sim\pi_k(\cdot\mid s)}\Bigg[D_\theta(s,v)\Bigg(
\nabla_\theta\log\pi_\theta(v\mid s)
-\mathbb E_{u\sim\pi_k(\cdot\mid s)}
\left[\nabla_\theta\log\pi_\theta(u\mid s)\right]\Bigg)\Bigg],\\
-\eta\nabla_\theta\mathcal L_{\mathrm{BPO}}(s)
&=\mathbb E_{v\sim\pi_k(\cdot\mid s)}
\left[\operatorname{stopgrad}\!\left(D_\theta(s,v)\right)
\nabla_\theta\log\pi_\theta(v\mid s)\right].
\end{align}
这就是前面删除第二项后的严格 policy-gradient 形式；被删除的项是
$\mathbb E_{u\sim\pi_k}[D_\theta(s,u)]
\mathbb E_{v\sim\pi_k}[\nabla_\theta\log\pi_\theta(v\mid s)]=0$，而不是把
后一个 score 期望本身当成零。

最后考虑 MDPO。记
\[
F_k(s,v)=\log\frac{T(v\mid s)}{\pi_k(v\mid s)},
\qquad
U_\theta(s,v)=\log\frac{\pi_\theta(v\mid s)}{\pi_k(v\mid s)}.
\]
对式~\eqref{eq:ch2-opd-full-mdpo-simplified} 求导时，既要微分外层测度
$\pi_\theta$，也要微分 $U_\theta$：
\begin{align}
\nabla_\theta\mathcal J_{\mathrm{MDPO}}(s)
&=\mathbb E_{v\sim\pi_\theta(\cdot\mid s)}
\Bigg[\left(F_k(s,v)-\frac1\eta U_\theta(s,v)-\frac1\eta\right)
\nabla_\theta\log\pi_\theta(v\mid s)\Bigg]\\
&=\mathbb E_{v\sim\pi_\theta(\cdot\mid s)}
\left[\operatorname{stopgrad}\!\left(F_k(s,v)-\frac1\eta U_\theta(s,v)\right)
\nabla_\theta\log\pi_\theta(v\mid s)\right],
\end{align}
因为 $\mathbb E_{v\sim\pi_\theta}[\nabla_\theta\log\pi_\theta(v\mid s)]=0$。因此 MDPO 是以当前 student
$\pi_\theta$ 为采样测度的 policy gradient，而 BPO/PMD-Mean 是以冻结 anchor
$\pi_k$ 为测度、通过 centered residual 形成的 policy gradient。

若 sampled-token 路径估计 $\mathcal K_k^\theta$，从 $a\sim\pi_k$ 采样时可以使用
\begin{equation}
\widehat{\mathcal K}_{k}^{\theta}(s,a)
=\frac{\pi_k(a\mid s)}{\operatorname{stopgrad}(\pi_k(a\mid s))}
\left[\operatorname{stopgrad}(\log\pi_k(a\mid s))
-\log\pi_\theta(a\mid s)\right].
\end{equation}
它在数值上是 $\log\pi_k(a\mid s)-\log\pi_\theta(a\mid s)$，但保留了
$-\nabla_\theta\log\pi_\theta(a\mid s)$。若把整个
$\widehat{\mathcal K}_{k}^{\theta}$ 错误地 stop-gradient，PMD-Mean 和 BPO
都会退化成未中心化的梯度，且不再对应式~\eqref{eq:ch2-opd-grad-mean-residual}。

\subsection{有限词表与 sampled-token：KL 与 partition 估计}
\label{sec:ch2-opd-z-kl-estimation}

2602.04417 的 Top-$K$ 方法直接估计 KL：在 top-$K$ 头部计算精确项，在尾部用
采样项校正。它不直接给出本节的 $\log Z_k$，因为
\begin{equation}
Z_k(s)=\mathbb E_{a\sim\pi_k(\cdot\mid s)}
\left[\left(\frac{T(a\mid s)}{\pi_k(a\mid s)}\right)^\eta\right].
\label{eq:ch2-opd-z-importance}
\end{equation}
这一区别很重要：$\widehat Z_k$ 可以无偏，但 $\log\widehat Z_k$ 不能无偏。

在固定 prefix $s$ 上，从冻结 anchor 采样 $a_j\sim\pi_k(\cdot\mid s)$，令
$r_k(s,a)=\log T(a\mid s)-\log\pi_k(a\mid s)$。则两个中心满足
\begin{align}
\mathcal K_k^T(s)&=\mathbb E_{a\sim\pi_k}[-r_k(s,a)],\\
\mathcal K_{k,n}^\theta(s)&=\mathbb E_{a\sim\pi_k}
\left[\log\pi_k(a\mid s)-\log\pi_{\theta_{k,n}}(a\mid s)\right].
\end{align}
因此 $m$ 个独立 next-token 样本的均值分别是两个 KL 的无偏估计。
真实 rollout 中每个 prefix 通常只有一个后继 token；该 token 是无偏单点样本，
但不是 conditional KL 的精确值。若要降低方差，可以在已经算出 anchor logits
后额外抽取 $m-1$ 个 token，这只增加 categorical 抽样和 log-prob 记录，
无需生成额外长序列。不同 prefix 的 token 不能互相当作同一条件分布的样本。

其中 $\mathcal K_k^T$ 在整个 outer step 内固定，teacher/anchor log-prob 可以
缓存；$\mathcal K_{k,n}^\theta$ 必须在每个 inner step 重新计算。实现已经推导
好的 policy-gradient 时，KL 中心作为权重整体使用 stop-gradient；实现原始
Mean/BPO 平方损失时则必须保留其对 student 的导数。

对 $\mathcal K_{k,n}^\theta=\operatorname{KL}(\pi_k\Vert\pi_{\theta_{k,n}})$，
取 $S_k(s)=\operatorname{TopK}(\pi_k(\cdot\mid s))$，可构造
\begin{align}
\widehat{\mathcal K}_{k,n}^{\theta}(s)
={}&\sum_{v\in S_k(s)}\pi_k(v\mid s)
\log\frac{\pi_k(v\mid s)}{\pi_{\theta_{k,n}}(v\mid s)}\\
&+\frac1m\sum_{j=1}^{m}\mathbf 1\{a_j\notin S_k(s)\}
\log\frac{\pi_k(a_j\mid s)}{\pi_{\theta_{k,n}}(a_j\mid s)},
\qquad a_j\sim\pi_k(\cdot\mid s).
\label{eq:ch2-opd-topk-anchor-kl}
\end{align}
对固定 $S_k$，尾部项是精确尾部求和的无偏 Monte Carlo 估计，且其 autograd
梯度的期望等于真实 KL 梯度。这里的“梯度无偏”是指把该估计量直接放进
原始 KL 损失并对 student 反向传播时的性质；本章的最终实现采用已经解析化的
policy-gradient 形式，KL 估计值只作为数值权重，必须整体
$\operatorname{stopgrad}$，因此不需要再对这个估计量执行 autograd。
$K=0$ 是纯 sampled-token，$K=|\mathcal V|$ 是 exact KL；这两个端点是同一
个 top-$K$ 族的插值端点，不单独构成一个算法分支。给定头部后，条件方差为
\begin{equation}
\operatorname{Var}(\widehat{\mathcal K}_{k,n}^{\theta}\mid S_k)
=\frac1m\operatorname{Var}_{a\sim\pi_k}\left[
\mathbf 1\{a\notin S_k\}
\log\frac{\pi_k(a\mid s)}{\pi_{\theta_{k,n}}(a\mid s)}\right].
\end{equation}
增大 $K$ 通常降低尾部质量和方差，但实际大小取决于尾部 log-ratio 的二阶矩。
对 $\mathcal K_k^T$ 可使用同样的头尾分解；由于它不需要反向传播，anchor
样本均值往往已经足够。

为了把 2602.04417 中“截断偏差换取采样方差”的结论直接用于 OPD，下面给出
一个与 KL、score 加权梯度和 partition function 都兼容的统一形式。固定 prefix
$s$，令 $p(v)$ 为 rollout 测度（本章中通常是 $p=\pi_k(\cdot\mid s)$），
$h(v)$ 为待求期望的 integrand，定义
\begin{equation}
Q=\mathbb E_{v\sim p}[h(v)],\qquad
Q_K^{\mathrm{head}}=\sum_{v\in S_K}p(v)h(v),\qquad
\tau_K=\mathbb E_{v\sim p}[\mathbf 1\{v\notin S_K\}h(v)],
\label{eq:ch2-opd-generic-tail}
\end{equation}
其中 $S_K$ 是固定的 top-$K$ 集合。用 $m$ 个独立的 auxiliary token
$a_1,\ldots,a_m\sim p$ 做尾部校正时，定义
\begin{equation}
\widehat Q_{K,m}=Q_K^{\mathrm{head}}+
\frac1m\sum_{j=1}^{m}\mathbf 1\{a_j\notin S_K\}h(a_j).
\label{eq:ch2-opd-generic-topk}
\end{equation}

\begin{theorem}[Top-$K$ 截断与尾部校正的偏差--方差分解]
\label{thm:ch2-opd-topk-bias-variance}
在 $\mathbb E_p[h^2]<\infty$ 且 $S_K$ 与 auxiliary samples 独立时，纯头部估计
$\widehat Q_K^{\mathrm{tr}}=Q_K^{\mathrm{head}}$ 与尾部校正估计
$\widehat Q_{K,m}$ 分别满足
\begin{align}
\operatorname{Bias}(\widehat Q_K^{\mathrm{tr}})&=-\tau_K,&
\operatorname{Var}(\widehat Q_K^{\mathrm{tr}})&=0,&
\operatorname{MSE}(\widehat Q_K^{\mathrm{tr}})&=\tau_K^2,\\
\operatorname{Bias}(\widehat Q_{K,m})&=0,&
\operatorname{Var}(\widehat Q_{K,m})&=\frac{V_K}{m},&
\operatorname{MSE}(\widehat Q_{K,m})&=\frac{V_K}{m},
\end{align}
其中
\begin{equation}
V_K=\mathbb E_p\left[\mathbf 1\{v\notin S_K\}h(v)^2\right]-\tau_K^2.
\label{eq:ch2-opd-tail-variance}
\end{equation}
如果 $S_K\subseteq S_{K+1}$，则尾部二阶矩
$A_K=\mathbb E_p[\mathbf 1\{v\notin S_K\}h(v)^2]$ 单调不增，且
$V_K\le A_K$。尾部均值 $\tau_K$ 在 $h\ge0$ 时单调不增；对符号不定的
$h$，它本身不必单调，但有
\begin{equation}
|\tau_K|\le \sqrt{p(S_K^c)A_K}.
\label{eq:ch2-opd-tail-cauchy}
\end{equation}
\end{theorem}

\begin{proof}
记 $X_j=\mathbf 1\{a_j\notin S_K\}h(a_j)$。由 $a_j\sim p$，有
$\mathbb E[X_j]=\tau_K$，且 $\operatorname{Var}(X_j)=V_K$。因此
\begin{equation}
\mathbb E[\widehat Q_{K,m}]
 =Q_K^{\mathrm{head}}+\tau_K=Q,
\qquad
\operatorname{Var}(\widehat Q_{K,m})
 =\operatorname{Var}\left(\frac1m\sum_{j=1}^mX_j\right)
 =\frac{V_K}{m},
\end{equation}
其中最后一步使用了独立性。纯头部估计没有随机项，且
$Q_K^{\mathrm{head}}-Q=-\tau_K$，从而得到前三个等式。若集合嵌套，
$\mathbf 1\{v\notin S_{K+1}\}\le\mathbf 1\{v\notin S_K\}$，故对非负
$h^2$ 有 $A_{K+1}\le A_K$。当 $h\ge0$ 时，尾部指示函数逐点下降，
所以 $0\le\tau_{K+1}\le\tau_K$。一般情形下对尾部随机变量应用
Cauchy--Schwarz 不等式即可得到式~\eqref{eq:ch2-opd-tail-cauchy}。
\end{proof}

定理~\ref{thm:ch2-opd-topk-bias-variance} 有两个直接含义。对
$h_{\mathrm{KL}}(v)=\log\frac{\pi_k(v\mid s)}{\pi_{\theta}(v\mid s)}$
或 $h_T(v)=\log\frac{\pi_k(v\mid s)}{T(v\mid s)}$，$h$ 可能正也可能负，
因此不能声称“增大 $K$ 必然单调减小截断 bias”；能够保证的是尾部二阶矩和
上界下降。对 $h_Z(v)=\exp(\eta r_k(s,v))\ge0$，尾部均值本身单调下降，
所以纯截断的 $Z_k$ 估计具有严格的单侧负偏差，而且随 $K$ 减小。
这正是 2602.04417 中 truncated KL 与 tail-corrected Top-$K$ KL 的数学区别。

如果最终对象是 score 加权的 policy gradient，令
$g_\theta(v)=\nabla_\theta\log\pi_\theta(v\mid s)$，把
$h(v)$ 换成向量 $h(v)g_\theta(v)$，定理逐坐标成立，并得到
\begin{equation}
\operatorname{Bias}(\widehat G_K^{\mathrm{tr}})
 =-\mathbb E_p[\mathbf 1\{v\notin S_K\}h(v)g_\theta(v)],
\qquad
\operatorname{Cov}(\widehat G_{K,m})
 =\frac1m\operatorname{Cov}_p(\mathbf 1\{v\notin S_K\}h(v)g_\theta(v)).
\label{eq:ch2-opd-gradient-topk-mse}
\end{equation}
这描述的是“估计最终梯度权重”的误差，而不是对 detached 估计量求导。
本章 Algorithm~5 采用后者：先解析得到 score-function 公式，再把
$\widehat h$ 整体 stop-gradient。因此实现不需要保存估计量的 autograd 图；
只需控制式~\eqref{eq:ch2-opd-gradient-topk-mse} 中的数值 bias 和 variance。

对 $Z_k$，头尾分解应先估计 $Z_k$ 再取对数。令
$S_Z(s)$ 为 $((1-\eta)\log\pi_k+\eta\log T)$ 的 top-$K$ 索引，则
\begin{align}
\widehat Z_k(s)
={}&\sum_{v\in S_Z(s)}\pi_k(v\mid s)^{1-\eta}T(v\mid s)^\eta\\
&+\frac1m\sum_{j=1}^{m}\mathbf 1\{a_j\notin S_Z(s)\}
\left(\frac{T(a_j\mid s)}{\pi_k(a_j\mid s)}\right)^\eta,
\qquad a_j\sim\pi_k(\cdot\mid s).
\label{eq:ch2-opd-topk-z}
\end{align}
式~\eqref{eq:ch2-opd-topk-z} 对 $Z_k$ 无偏，但 Jensen 不等式给出
\begin{equation}
\mathbb E[\log\widehat Z_k]\leq\log Z_k,
\qquad
\mathbb E[\log\widehat Z_k]-\log Z_k
\approx-\frac{\operatorname{Var}(\widehat Z_k)}{2Z_k^2}.
\end{equation}
尾部重要性比率具有长尾时，单 token plug-in 估计尤其不稳定；裁剪比率可降方差，
但会引入额外偏差。$\eta=1$ 时 $Z_k=1$，所以 $\log Z_k=0$ 精确已知；
$\eta=0$ 时同样如此。$0<\eta<1$ 时，若不需要 PMD-Part，PMD-Mean/BPO
可以绕开这项 log-of-mean 偏差。

\begin{theorem}[Partition function 的 log-of-mean 偏差]
\label{thm:ch2-opd-logz-delta}
设 $\widehat Z_{K,m}$ 是式~\eqref{eq:ch2-opd-topk-z} 的尾部校正估计，
$\mathbb E[\widehat Z_{K,m}]=Z_k>0$，且尾部三阶绝对矩有限。记
$V_{K}=\operatorname{Var}(\mathbf 1\{a\notin S_Z\}
(T(a\mid s)/\pi_k(a\mid s))^\eta)$。当 $m$ 足够大时，
\begin{align}
\mathbb E[\log\widehat Z_{K,m}]-\log Z_k
&=-\frac{V_K}{2mZ_k^2}+O(m^{-3/2}),\\
\operatorname{Var}(\log\widehat Z_{K,m})
&=\frac{V_K}{mZ_k^2}+O(m^{-3/2}).
\end{align}
若改用纯头部估计 $Z_K^{\mathrm{tr}}=\sum_{v\in S_Z}
\pi_k(v)^{1-\eta}T(v)^\eta$，则
\begin{equation}
Z_K^{\mathrm{tr}}-Z_k=-\tau_K^Z\le0,
\qquad
\log Z_K^{\mathrm{tr}}-\log Z_k\le0,
\end{equation}
其中 $\tau_K^Z\ge0$ 是式~\eqref{eq:ch2-opd-generic-tail} 对
$h_Z$ 的尾部均值。因而纯截断的 log-partition bias 是确定性的单侧偏差，
而尾部校正的 plug-in log 估计是由方差诱导的负偏差。
\end{theorem}

\begin{proof}
令 $\Delta=\widehat Z_{K,m}-Z_k$。由无偏性，$\mathbb E[\Delta]=0$，且
$\mathbb E[\Delta^2]=V_K/m$。在 $Z_k$ 邻域对
$f(z)=\log z$ 作三阶 Taylor 展开：
\begin{equation}
\log(Z_k+\Delta)=\log Z_k+\frac{\Delta}{Z_k}
 -\frac{\Delta^2}{2Z_k^2}+R_3,
\end{equation}
其中在三阶矩有界时 $\mathbb E|R_3|=O(m^{-3/2})$。取期望并使用
$\mathbb E[\Delta]=0$ 得到第一式。对同一展开取方差，首阶随机项为
$\Delta/Z_k$，故
$\operatorname{Var}(\log\widehat Z_{K,m})=V_K/(mZ_k^2)+O(m^{-3/2})$。
最后，$h_Z\ge0$ 给出 $\tau_K^Z\ge0$，所以
$Z_K^{\mathrm{tr}}\le Z_k$；由于 log 单调递增，第二个不等式随之成立。
\end{proof}

定理~\ref{thm:ch2-opd-logz-delta} 说明了为什么不能把“无偏的
$\widehat Z_k$”直接等同于“无偏的 $\log Z_k$”。在 PMD-Part 中，若使用
plug-in 估计，$\log Z_k$ 的误差会进入每个 token 的 residual；在完整词表
路径上 score 均值可精确消去该常数项，在 sampled-token 路径上只能依赖有限样本
近似消去。因此 PMD-Part 对 $m$、$K$ 和尾部比率的要求严格于 PMD-Mean/BPO。
当 $m$ 个样本独立且把同一 prefix 的梯度平均时，partition 诱导的梯度方差
按 $1/m$ 下降；当外层有 $B$ 个相互独立 prefix 时，条件于 prefix 的这部分
方差进一步按 $1/B$ 下降，但 prefix coverage 误差不随 $m$ 消失。

\subsubsection*{偏差--方差定理：从 KL 数值到最终 policy gradient}
对一个 prefix 的任意估计器 $\widehat q$，最终关心的是带 score 的量
$\widehat g=\widehat q\,\nabla_\theta\log\pi_\theta$。其均方误差分解为
\begin{equation}
\mathbb E\lVert\widehat g-g\rVert^2
=\lVert\mathbb E[\widehat g]-g\rVert^2
+\operatorname{tr}\operatorname{Cov}(\widehat g).
\end{equation}
因此，单独比较 $\widehat{\mathcal K}$ 的数值方差并不足够：一个有轻微
截断偏差但尾部方差极低的估计器，可能比无偏但重尾的 sampled estimator
产生更小的梯度 MSE。2602.04417 的 Top-$K$ 结论正是这个取舍：加入尾部校正
保持无偏，代价是尾部采样方差；只保留头部则引入可控的截断偏差。

对 OPD 的 PMD-Part，$\log Z_k$ 的误差直接进入 residual，而 residual 又
乘以 score。若使用 $\widehat c_k=\log\widehat Z_k$，令
$\delta_Z=\widehat c_k-\log Z_k$，则单 prefix 的梯度误差包含
\begin{equation}
\Delta g_Z
=\delta_Z\,\mathbb E_{v\sim\pi_k}
\left[\nabla_\theta\log\pi_\theta(v\mid s)\right]
\end{equation}
以及 $\delta_Z$ 与 sampled score 的协方差项。完整词表下 score 均值可以被
精确求和，前一项严格为零；sampled-token 下它只是在期望上为零，有限 batch
会把 $\log Z_k$ 的标量误差放大成随机梯度噪声。故 PMD-Part 对 partition
估计的要求比 BPO/PMD-Mean 更严格。

对同一 prefix 同时估计中心和 policy-gradient 权重时，还要考虑协方差：若
$\widehat c$ 与 score 使用同一枚 token，$\mathbb E[\widehat c\,\widehat g]$
一般不等于 $\mathbb E[\widehat c]\mathbb E[\widehat g]$。采用独立的 auxiliary
next-token 样本估计中心，或使用 leave-one-out center，可以消除这类自耦合偏差；
代价是额外的 token 抽样和少量 log-prob 存储。跨 prefix 聚合只能降低总体
均值方差，不能把不同条件分布下的样本变成同一 conditional center。

最后，prefix 覆盖误差与 token 估计误差应分开记录。即使 token-level estimator
在给定 $s$ 上无偏，rollout 只访问由 $\pi_k$ 诱导的有限 prefix 多重集，仍有
$\widehat d_k-d_k$ 的 outer 测度误差。增加 Top-$K$ 或 auxiliary samples 只
降低条件 token 方差，不能修复 prefix coverage；增加 rollout 数量或重新采样
outer batch 才能降低这一项。实验日志应至少分别保存：prefix 数量与重复率、
KL/partition 的 batch 均值和方差、尾部采样次数、梯度范数与峰值显存。

缓存上，CPU/NVMe 只需保存 token IDs、mask、response mask、anchor/teacher
的 action log-prob 和可选的 top-$K$ 摘要；student log-prob 每个 inner step
重算。保存完整 logits 的代价是 $O(BH|\mathcal V|)$，保存 top-$K$ 摘要是
$O(BHK)$。以 $|\mathcal V|=150{,}000$、FP16 为例，一个 prefix-token 的
完整 logits 约 $300$ kB，8192 token 单序列约 $2.46$ GB；因此不能把全词表
logits 放进 replay buffer。GPU 只保留当前 microbatch 的 student 激活和当前
token 块 logits，并在块处理后释放。当前 OPD 最稳妥的首版路径是缓存 anchor/
teacher action log-prob，用额外 anchor next-token 样本估计两个 KL 中心；
PMD-Part 只在 $\eta=1$ 或已验证 $\log Z_k$ 误差足够小时启用，再用 top-$K$
实验比较中心方差、$\log Z_k$ 偏差和峰值显存。

\begin{algorithm}[ht]
\caption{OPD full-vocabulary outer--inner optimization\\{\small\textit{Semantics: 每个 outer step 用 $\pi_k$ rollout，固定 anchor 做 $N$ 个 inner policy-gradient step，再同步得到 $\pi_{k+1}$}}}
\begin{algorithmic}[1]
\Require prompt distribution $\mathcal D$, teacher $T$, initial student $\pi_0=\pi_{\theta_{0,0}}$, outer steps $K$, rollout batch size $B$, horizon $H$, method $M$, $\eta$, inner steps $N$, inner learning rate schedule $\{\alpha_{k,n}\}$, sync mode $\mathrm{Sync}\in\{\mathrm{hard},\mathrm{soft}\}$, soft-sync weight $\rho\in(0,1]$
\Ensure final policy $\pi_K$ and parameters $\theta_K$
\For{$k=0,\ldots,K-1$}
    \State $\mathcal B_k\gets\varnothing$ \Comment{$\mathcal B_k$ is a multiset; repeated prefixes keep their multiplicity}
    \For{$i=1,\ldots,B$}
        \State sample prompt $x_i\sim\mathcal D$
        \State generate $y_i=(a_{i,1},\ldots,a_{i,\tau_i})$ with $a_{i,t}\sim\pi_k(\cdot\mid s_{i,t})$, $s_{i,t}=(x_i,a_{i,<t})$
        \State $\mathcal B_k\gets\mathcal B_k\mathbin{\uplus}\{s_{i,t}:1\le t\le\tau_i\wedge H\}$
    \EndFor
    \State define $\widehat d_k$ as the empirical distribution assigning each occurrence in $\mathcal B_k$ equal mass
    \State $\theta_{k,0}\gets\theta_k$ \Comment{initialize the inner student from the current outer policy}
    \For{$n=0,\ldots,N-1$}
        \State query $\log T(v\mid s)$, $\log\pi_k(v\mid s)$ and $\log\pi_{\theta_{k,n}}(v\mid s)$ for every occurrence $s\in\mathcal B_k$ and $v\in\mathcal V$
        \State $\log Z_k(s)\gets\log\left(\sum_{v\in\mathcal V}\exp\left((1-\eta)\log\pi_k(v\mid s)+\eta\log T(v\mid s)\right)\right)$
        \State $\mathcal K_k^T(s)\gets\operatorname{KL}(\pi_k(\cdot\mid s)\Vert T(\cdot\mid s))$
        \State $\mathcal K_{k,n}^\theta(s)\gets\operatorname{KL}(\pi_k(\cdot\mid s)\Vert\pi_{\theta_{k,n}}(\cdot\mid s))$
        \If{$M=\mathrm{Part}$}
            \State $G_{k,n}\gets\mathbb E_{s\sim\widehat d_k}\mathbb E_{v\sim\pi_k}\left[\operatorname{stopgrad}\!\left(\log\frac{\pi_{\theta_{k,n}}}{\pi_k}-\eta\log\frac{T}{\pi_k}+\log Z_k\right)\nabla_{\theta_{k,n}}\log\pi_{\theta_{k,n}}\right]$
        \ElsIf{$M=\mathrm{Mean}$}
            \State $G_{k,n}\gets\mathbb E_{s\sim\widehat d_k}\mathbb E_{v\sim\pi_k}\left[\operatorname{stopgrad}\!\left(\log\frac{\pi_{\theta_{k,n}}}{\pi_k}+\mathcal K_{k,n}^\theta-\eta\left(\log\frac{T}{\pi_k}+\mathcal K_k^T\right)\right)\nabla_{\theta_{k,n}}\log\pi_{\theta_{k,n}}\right]$
        \ElsIf{$M=\mathrm{BPO}$}
            \State $G_{k,n}\gets-\eta^{-1}\mathbb E_{s\sim\widehat d_k}\mathbb E_{v\sim\pi_k}\left[\operatorname{stopgrad}\!\left(\eta\left(\log\frac{T}{\pi_k}+\mathcal K_k^T\right)-\left(\log\frac{\pi_{\theta_{k,n}}}{\pi_k}+\mathcal K_{k,n}^\theta\right)\right)\nabla_{\theta_{k,n}}\log\pi_{\theta_{k,n}}\right]$
        \ElsIf{$M=\mathrm{MDPO}$}
            \State $G_{k,n}\gets-\mathbb E_{s\sim\widehat d_k}\mathbb E_{v\sim\pi_{\theta_{k,n}}}\left[\operatorname{stopgrad}\!\left(\log\frac{T}{\pi_k}-\eta^{-1}\log\frac{\pi_{\theta_{k,n}}}{\pi_k}\right)\nabla_{\theta_{k,n}}\log\pi_{\theta_{k,n}}\right]$
        \EndIf
        \State $\theta_{k,n+1}\gets\operatorname{OptimizerStep}(\theta_{k,n},G_{k,n};\alpha_{k,n})$ \Comment{expectations are implemented by exact sums or MC averages}
    \EndFor
    \If{$\mathrm{Sync}=\mathrm{hard}$}
        \State $\theta_{k+1}\gets\theta_{k,N}$ and $\pi_{k+1}\gets\pi_{\theta_{k+1}}$
    \Else
        \State $\pi_{k+1}(\cdot\mid s)\gets(1-\rho)\pi_k(\cdot\mid s)+\rho\pi_{\theta_{k,N}}(\cdot\mid s)$ for every prefix $s$
        \State choose $\theta_{k+1}$ to represent $\pi_{k+1}$ \Comment{in parameter-EMA implementations, set $\theta_{k+1}=(1-\rho)\theta_k+\rho\theta_{k,N}$}
    \EndIf
\EndFor
\State \Return $(\pi_K,\theta_K)$
\end{algorithmic}
\end{algorithm}

算法中的同步发生在整个 inner loop 结束之后。因而在固定 $k$ 的所有
$N$ 次更新中，rollout 测度、anchor ratio $\pi_k/T$ 以及所有 mirror
target 都保持不变；只有 student 参数 $\theta_{k,n}$ 变化。hard sync
令下一轮 anchor 完全等于 inner 末端 student，soft sync 则令下一轮
anchor 是旧 anchor 与 inner 末端 student 的凸组合。完成同步后丢弃
$\widehat d_k$，使用新的 $\pi_{k+1}$ 重新 rollout，而不是继续复用旧
rollout。若 inner loop 没有收敛，$\pi_{\theta_{k,N}}$ 只是当前固定
mirror 子问题的有限步近似解；outer 动力学仍由上述同步映射严格定义。

\subsubsection*{只有 sampled next token：无完整词表时的估计路径}

如果每个 prefix 只有 rollout 中的下一个 token $a\sim\pi_k(\cdot\mid s)$，仍然可以无偏观测 $r_k(s,a)$，但不能从一个样本无偏得到 $\log Z_k(s)$：先估计 $Z_k$ 再取对数会产生 plug-in bias。因此 sampled-token 路径不应把 PMD-Part 的 exact $\log Z_k$ 假装已知。可行选择如下。

第一，使用 BPO 或 PMD-Mean，以同一个 prefix 的多条 sampled actions 估计 conditional center；它们不需要显式的 $\log Z_k$。第二，使用 MDPO 的 action-sampled objective，并用 K4/K3$^{++}$ 估计 KL 项；这要求保留 student/anchor 的 sampled log-prob。第三，如果只有一条 action，却仍想做 PMD-Part，只能使用外部 baseline、跨 batch center 或 self-normalized $\widehat{\log Z}$，并在理论和日志中标记其 bias。

对单个 sampled token，PMD-Part 的可实现 loss 是带 baseline 的近似
\[
\widehat\ell_{\mathrm{Part}}=\frac12\bigl[\eta(r_k(s,a)-b(s))-U_\theta(s,a)\bigr]^2,
\]
它只有在 $b(s)=\log Z_k(s)/\eta$ 时才等于 exact residual；任意其他 baseline 会留下一个 prefix-dependent constant，因而改变 PMD-Part 的有限参数投影。BPO 的中心化 residual 则可用同 prefix 的 $K$ 个 sampled tokens 构造，$K=1$ 时必须使用跨样本或 frozen center。

\begin{algorithm}[ht]
\caption{OPD sampled-token inner update\\{\small\textit{Semantics: 每个 prefix 只访问 rollout 的下一个 token}}}
\begin{algorithmic}[1]
\Require frozen anchor $\pi_k$, teacher $T$, rollout pairs $(s_i,a_i)$, method $M$, center size $K$, $\eta$
\Ensure updated student parameters $\theta$
\For{each prefix $s_i$}
    \State read $\log T(a_i\mid s_i)$ and $\log\pi_k(a_i\mid s_i)$
    \State $r_i\gets\log T(a_i\mid s_i)-\log\pi_k(a_i\mid s_i)$
    \If{$M=\mathrm{BPO}$ or $M=\mathrm{Mean}$}
        \State obtain $K$ actions from independent same-prefix rollouts, or use a frozen cross-batch center
        \State compute leave-one-out center of $r$ and $U_\theta$; keep the center's student gradient
        \State update with the corresponding centered residual loss
    \ElsIf{$M=\mathrm{MDPO}$}
        \State form $R_i=\exp(\log\pi_\theta(a_i\mid s_i)-\operatorname{stopgrad}(\log\pi_k(a_i\mid s_i)))$
        \State use $R_i r_i-\eta^{-1}\widehat{\mathrm{KL}}_i$ with K4/K3$^{++}$ and retain correct autograd paths
    \ElsIf{$M=\mathrm{Part}$}
        \State estimate or load a prefix baseline $\widehat{\log Z_i}$; record its bias source
        \State update $\frac12[\eta r_i-\widehat{\log Z_i}-U_\theta(s_i,a_i)]^2$
    \EndIf
\EndFor
\State update $\theta$ while keeping $\pi_k$ and $T$ frozen
\State \Return $\theta$
\end{algorithmic}
\end{algorithm}

完整 vocabulary 与 sampled-token 两条路径的区别是信息论上的：前者能精确知道整个 categorical mirror target，后者只能无偏知道在 anchor 下抽到的一个 action 的 reward。两者在无限 samples、无函数类误差且 center/partition 估计一致时收敛到同一个 token fixed point；有限 batch 下，full-vocabulary 的主要误差是参数优化和 prefix coverage，sampled-token 额外承担 reward、center、KL ratio 和 prefix reuse 的估计方差。
